\Chapter{109}

\Verse{1}
{
    \zR{话说宝钗叫袭人问出原故，恐宝玉悲伤成疾，便将黛玉临死的话与袭人假作闲 谈，说是：“人在世上，有意有情，到了死后，各自干各自的去了，并不是生前那
        样的人死后还是那样。}
}
{
    \eR{[English source not present in the checked-in Bencraft/Joly files.]}
}
{
}

\Verse{2}
{
    \zR{活人虽有痴心，死的竟不知道。}
}
{
    \eR{[English source not present in the checked-in Bencraft/Joly files.]}
}
{
}

\Verse{3}
{
    \zR{况且林姑娘既说仙去，他看 凡人是个不堪的浊物，那里还肯混在世上？}
}
{
    \eR{[English source not present in the checked-in Bencraft/Joly files.]}
}
{
}

\Verse{4}
{
    \zR{只是人自己疑心，所以招出些邪魔外祟 来缠扰。”}
}
{
    \eR{[English source not present in the checked-in Bencraft/Joly files.]}
}
{
}

\Verse{5}
{
    \zR{宝钗虽是与袭人说话，原说给宝玉听的。}
}
{
    \eR{[English source not present in the checked-in Bencraft/Joly files.]}
}
{
}

\Verse{6}
{
    \zR{袭人会意，也说是：“没有的 事。}
}
{
    \eR{[English source not present in the checked-in Bencraft/Joly files.]}
}
{
}

\Verse{7}
{
    \zR{若说林姑娘的魂灵儿还在园里，我们也算相好，怎么没有梦见过一次？”}
}
{
    \eR{[English source not present in the checked-in Bencraft/Joly files.]}
}
{
}

\Verse{8}
{
    \zR{宝玉 在外面听着，细细的想道：“果然也奇。}
}
{
    \eR{[English source not present in the checked-in Bencraft/Joly files.]}
}
{
}

\Verse{9}
{
    \zR{我知道林妹妹死了，那一日不想几遍，怎 么从没梦见？}
}
{
    \eR{[English source not present in the checked-in Bencraft/Joly files.]}
}
{
}

\Verse{10}
{
    \zR{想必他到天上去了，瞧我这凡夫俗子不能交通神明，所以梦都没有一 个儿。}
}
{
    \eR{[English source not present in the checked-in Bencraft/Joly files.]}
}
{
}

\Verse{11}
{
    \zR{我如今就在外间睡，或者我从园里回来，他知道我的心，肯与我梦里一见。}
}
{
    \eR{[English source not present in the checked-in Bencraft/Joly files.]}
}
{
}

\Verse{12}
{
    \zR{我必要问他实在那里去了，我也时常祭奠。}
}
{
    \eR{[English source not present in the checked-in Bencraft/Joly files.]}
}
{
}

\Verse{13}
{
    \zR{若是果然不理我这浊物，竟无一梦，我 便也不想他了。”}
}
{
    \eR{[English source not present in the checked-in Bencraft/Joly files.]}
}
{
}

\Verse{14}
{
    \zR{主意已定，便说：“我今夜就在外间睡，你们也不用管我。”}
}
{
    \eR{[English source not present in the checked-in Bencraft/Joly files.]}
}
{
}

\Verse{15}
{
    \zR{宝 钗也不强他，只说：“你不用胡思乱想。}
}
{
    \eR{[English source not present in the checked-in Bencraft/Joly files.]}
}
{
}

\Verse{16}
{
    \zR{你没瞧见太太因你园里去了，急的话都说 不出来？}
}
{
    \eR{[English source not present in the checked-in Bencraft/Joly files.]}
}
{
}

\Verse{17}
{
    \zR{你这会子还不保养身子，倘或老太太知道了，又说我们不用心。”}
}
{
    \eR{[English source not present in the checked-in Bencraft/Joly files.]}
}
{
}

\Verse{18}
{
    \zR{宝玉道： “白这么说罢咧，我坐一会子就进来。}
}
{
    \eR{[English source not present in the checked-in Bencraft/Joly files.]}
}
{
}

\Verse{19}
{
    \zR{你也乏了，先睡罢。”}
}
{
    \eR{[English source not present in the checked-in Bencraft/Joly files.]}
}
{
}

\Verse{20}
{
    \zR{宝钗料他必进来的， 假意说道：“我睡了，叫袭姑娘伺候你罢。”}
}
{
    \eR{[English source not present in the checked-in Bencraft/Joly files.]}
}
{
}

\Verse{21}
{
    \zR{宝玉听了，正合机宜。}
}
{
    \eR{[English source not present in the checked-in Bencraft/Joly files.]}
}
{
}

\Verse{22}
{
    \zR{等宝钗睡下，他便叫袭人麝月另铺设下一副被褥，常叫 人进来瞧二奶奶睡着了没有。}
}
{
    \eR{[English source not present in the checked-in Bencraft/Joly files.]}
}
{
}

\Verse{23}
{
    \zR{宝钗故意装睡，也是一夜不宁。}
}
{
    \eR{[English source not present in the checked-in Bencraft/Joly files.]}
}
{
}

\Verse{24}
{
    \zR{那宝玉只当宝钗睡着， 便与袭人道：“你们各自睡罢，我又不伤感。}
}
{
    \eR{[English source not present in the checked-in Bencraft/Joly files.]}
}
{
}

\Verse{25}
{
    \zR{你若不信，你就伏侍我睡了再进去， 只要不惊动我就是了。”}
}
{
    \eR{[English source not present in the checked-in Bencraft/Joly files.]}
}
{
}

\Verse{26}
{
    \zR{袭人果然伏侍他睡下，预备下了茶水，关好了门，进里间 去照应了一回，各自假寐，等着宝玉若有动静再出来。}
}
{
    \eR{[English source not present in the checked-in Bencraft/Joly files.]}
}
{
}

\Verse{27}
{
    \zR{宝玉见袭人进去了，便将坐 更的两个婆子支到外头。}
}
{
    \eR{[English source not present in the checked-in Bencraft/Joly files.]}
}
{
}

\Verse{28}
{
    \zR{他轻轻的坐起来，暗暗的祝赞了几句，方才睡下。}
}
{
    \eR{[English source not present in the checked-in Bencraft/Joly files.]}
}
{
}

\Verse{29}
{
    \zR{起初再 睡不着，以后把心一静，谁知竟睡着了，却倒一夜安眠。}
}
{
    \eR{[English source not present in the checked-in Bencraft/Joly files.]}
}
{
}

\Verse{30}
{
    \zR{直到天亮，方才醒来，拭 了拭眼，坐着想了一回，并无有梦。}
}
{
    \eR{[English source not present in the checked-in Bencraft/Joly files.]}
}
{
}

\Verse{31}
{
    \zR{便叹口气道：“正是‘悠悠生死别经年，魂魄 不曾来入梦’！”}
}
{
    \eR{[English source not present in the checked-in Bencraft/Joly files.]}
}
{
}

\Verse{32}
{
    \zR{宝钗反是一夜没有睡着，听见宝玉在外边念这两句，便接口道： “这话你说莽撞了。}
}
{
    \eR{[English source not present in the checked-in Bencraft/Joly files.]}
}
{
}

\Verse{33}
{
    \zR{若林妹妹在时，又该生气了。”}
}
{
    \eR{[English source not present in the checked-in Bencraft/Joly files.]}
}
{
}

\Verse{34}
{
    \zR{宝玉听了，自觉不好意思，只 得起来，搭讪着进里间来，说：“我原要进来，不知怎么一个盹儿就打着了。”}
}
{
    \eR{[English source not present in the checked-in Bencraft/Joly files.]}
}
{
}

\Verse{35}
{
    \zR{宝 钗道：“你进来不进来，与我什么相干？”}
}
{
    \eR{[English source not present in the checked-in Bencraft/Joly files.]}
}
{
}

\Verse{36}
{
    \zR{袭人也本没有睡，听见他们两个说话，即忙上来倒茶。}
}
{
    \eR{[English source not present in the checked-in Bencraft/Joly files.]}
}
{
}

\Verse{37}
{
    \zR{只见老太太那边打发小 丫头来问：“宝二爷昨夜睡的安顿么？}
}
{
    \eR{[English source not present in the checked-in Bencraft/Joly files.]}
}
{
}

\Verse{38}
{
    \zR{若安顿，早早的同二奶奶梳洗了就过去。”}
}
{
    \eR{[English source not present in the checked-in Bencraft/Joly files.]}
}
{
}

\Verse{39}
{
    \zR{袭人道：“你去回老太太，说：‘宝玉昨夜很安顿，回来就过来。’”}
}
{
    \eR{[English source not present in the checked-in Bencraft/Joly files.]}
}
{
}

\Verse{40}
{
    \zR{小丫头去了。}
}
{
    \eR{[English source not present in the checked-in Bencraft/Joly files.]}
}
{
}

\Verse{41}
{
    \zR{宝钗连忙梳洗了，莺儿袭人等跟着，先到贾母那里行了礼。}
}
{
    \eR{[English source not present in the checked-in Bencraft/Joly files.]}
}
{
}

\Verse{42}
{
    \zR{便到王夫人那边起，至 凤姐，都让过了。}
}
{
    \eR{[English source not present in the checked-in Bencraft/Joly files.]}
}
{
}

\Verse{43}
{
    \zR{仍到贾母处，见他母亲也过来了。}
}
{
    \eR{[English source not present in the checked-in Bencraft/Joly files.]}
}
{
}

\Verse{44}
{
    \zR{大家问起：“宝玉晚上好么？”}
}
{
    \eR{[English source not present in the checked-in Bencraft/Joly files.]}
}
{
}

\Verse{45}
{
    \zR{宝钗便说：“回去就睡了，没有什么。”}
}
{
    \eR{[English source not present in the checked-in Bencraft/Joly files.]}
}
{
}

\Verse{46}
{
    \zR{众人放心，又说些闲话。}
}
{
    \eR{[English source not present in the checked-in Bencraft/Joly files.]}
}
{
}

\Verse{47}
{
    \zR{只见小丫头进来，说：“二姑奶奶要回去了。}
}
{
    \eR{[English source not present in the checked-in Bencraft/Joly files.]}
}
{
}

\Verse{48}
{
    \zR{听见说，孙姑爷那边人来，到大 太太那里说了些话，大太太叫人到四姑娘那边说，不必留了，让他去罢。}
}
{
    \eR{[English source not present in the checked-in Bencraft/Joly files.]}
}
{
}

\Verse{49}
{
    \zR{如今二姑 奶奶在大太太那边哭呢，大约就过来辞老太太。”}
}
{
    \eR{[English source not present in the checked-in Bencraft/Joly files.]}
}
{
}

\Verse{50}
{
    \zR{贾母众人听了，心中好不自在， 都说：“二姑娘这么一个人，为什么命里遭着这样的人!一辈子不能出头，这可怎 么好呢。”}
}
{
    \eR{[English source not present in the checked-in Bencraft/Joly files.]}
}
{
}

\Verse{51}
{
    \zR{说着，迎春进来，泪痕满面。}
}
{
    \eR{[English source not present in the checked-in Bencraft/Joly files.]}
}
{
}

\Verse{52}
{
    \zR{因是宝钗的好日子，只得含着泪，辞了众 人要回去。}
}
{
    \eR{[English source not present in the checked-in Bencraft/Joly files.]}
}
{
}

\Verse{53}
{
    \zR{贾母知道他的苦处，也不便强留，只说道：“你回去也罢了，但只不用 伤心。}
}
{
    \eR{[English source not present in the checked-in Bencraft/Joly files.]}
}
{
}

\Verse{54}
{
    \zR{碰着这样人也是没法儿的。}
}
{
    \eR{[English source not present in the checked-in Bencraft/Joly files.]}
}
{
}

\Verse{55}
{
    \zR{过几天我再打发人接你去罢。”}
}
{
    \eR{[English source not present in the checked-in Bencraft/Joly files.]}
}
{
}

\Verse{56}
{
    \zR{迎春道：“老太 太始终疼我，如今也疼不来了。}
}
{
    \eR{[English source not present in the checked-in Bencraft/Joly files.]}
}
{
}

\Verse{57}
{
    \zR{可怜我没有再来的时候儿了。”}
}
{
    \eR{[English source not present in the checked-in Bencraft/Joly files.]}
}
{
}

\Verse{58}
{
    \zR{说着，眼泪直流。}
}
{
    \eR{[English source not present in the checked-in Bencraft/Joly files.]}
}
{
}

\Verse{59}
{
    \zR{众人都劝道：“这有什么不能回来的呢？}
}
{
    \eR{[English source not present in the checked-in Bencraft/Joly files.]}
}
{
}

\Verse{60}
{
    \zR{比不得你三妹妹隔得远，要见面就难了。”}
}
{
    \eR{[English source not present in the checked-in Bencraft/Joly files.]}
}
{
}

\Verse{61}
{
    \zR{贾母等想起探春，不觉也大家落泪。}
}
{
    \eR{[English source not present in the checked-in Bencraft/Joly files.]}
}
{
}

\Verse{62}
{
    \zR{为是宝钗的生日，只得转悲作喜说：“这也不 难。}
}
{
    \eR{[English source not present in the checked-in Bencraft/Joly files.]}
}
{
}

\Verse{63}
{
    \zR{只要海疆平静，那边亲家调进京来，就见的着了。”}
}
{
    \eR{[English source not present in the checked-in Bencraft/Joly files.]}
}
{
}

\Verse{64}
{
    \zR{大家说：“可不是这么着 么？”}
}
{
    \eR{[English source not present in the checked-in Bencraft/Joly files.]}
}
{
}

\Verse{65}
{
    \zR{说着，迎春只得含悲而别。}
}
{
    \eR{[English source not present in the checked-in Bencraft/Joly files.]}
}
{
}

\Verse{66}
{
    \zR{大家送了出来，仍回贾母那里。}
}
{
    \eR{[English source not present in the checked-in Bencraft/Joly files.]}
}
{
}

\Verse{67}
{
    \zR{从早至暮，又闹 了一天，众人见贾母劳乏，各自散了。}
}
{
    \eR{[English source not present in the checked-in Bencraft/Joly files.]}
}
{
}

\Verse{68}
{
    \zR{独有薛姨妈辞了贾母，到宝钗那里，说道：“你哥哥是今年过了，直要等到皇 恩大赦的时候，减了等，才好赎罪。}
}
{
    \eR{[English source not present in the checked-in Bencraft/Joly files.]}
}
{
}

\Verse{69}
{
    \zR{这几年叫我孤苦伶仃，怎么处!我想要给你二 哥哥完婚，你想想好不好？”}
}
{
    \eR{[English source not present in the checked-in Bencraft/Joly files.]}
}
{
}

\Verse{70}
{
    \zR{宝钗道：“妈妈是因为大哥哥娶了亲，唬怕了的，所 以把二哥哥的事也疑惑起来。}
}
{
    \eR{[English source not present in the checked-in Bencraft/Joly files.]}
}
{
}

\Verse{71}
{
    \zR{据我说，很该办。}
}
{
    \eR{[English source not present in the checked-in Bencraft/Joly files.]}
}
{
}

\Verse{72}
{
    \zR{邢姑娘是妈妈知道的，如今在这里 也很苦。}
}
{
    \eR{[English source not present in the checked-in Bencraft/Joly files.]}
}
{
}

\Verse{73}
{
    \zR{娶了去，虽说咱们穷，究竟比他傍人门户好多着呢。”}
}
{
    \eR{[English source not present in the checked-in Bencraft/Joly files.]}
}
{
}

\Verse{74}
{
    \zR{薛姨妈道：“你得 便的时候，就去回明老太太，说我家没人，就要择日子了。”}
}
{
    \eR{[English source not present in the checked-in Bencraft/Joly files.]}
}
{
}

\Verse{75}
{
    \zR{宝钗道：“妈妈只管 和二哥哥商量，挑个好日子，过来和老太太、大太太说了，娶过去，就完了一宗事。}
}
{
    \eR{[English source not present in the checked-in Bencraft/Joly files.]}
}
{
}

\Verse{76}
{
    \zR{这里大太太也巴不得娶了去才好。”}
}
{
    \eR{[English source not present in the checked-in Bencraft/Joly files.]}
}
{
}

\Verse{77}
{
    \zR{薛姨妈道：“今日听见史姑娘也就回去了，老 太太心里要留你妹妹在这里住几天，所以他住下了。}
}
{
    \eR{[English source not present in the checked-in Bencraft/Joly files.]}
}
{
}

\Verse{78}
{
    \zR{我想他也是不定多早晚就走的 人了，你们姐妹们也多叙几天话儿。”}
}
{
    \eR{[English source not present in the checked-in Bencraft/Joly files.]}
}
{
}

\Verse{79}
{
    \zR{宝钗道：“正是呢。”}
}
{
    \eR{[English source not present in the checked-in Bencraft/Joly files.]}
}
{
}

\Verse{80}
{
    \zR{于是薛姨妈又坐了一 坐，出来辞了众人回去了。}
}
{
    \eR{[English source not present in the checked-in Bencraft/Joly files.]}
}
{
}

\Verse{81}
{
    \zR{却说宝玉晚间归房，因想：“昨夜黛玉竟不入梦，或者他已经成仙，所以不肯 来见我这种浊人，也是有的；不然，就是我的性儿太急了，也未可知。”}
}
{
    \eR{[English source not present in the checked-in Bencraft/Joly files.]}
}
{
}

\Verse{82}
{
    \zR{便想了个 主意，向宝钗说道：“我昨夜偶然在外头睡着，似乎比在屋里睡的安稳些，今日起 来，心里也觉清净。}
}
{
    \eR{[English source not present in the checked-in Bencraft/Joly files.]}
}
{
}

\Verse{83}
{
    \zR{我的意思，还要在外头睡两夜，只怕你们又来拦我。”}
}
{
    \eR{[English source not present in the checked-in Bencraft/Joly files.]}
}
{
}

\Verse{84}
{
    \zR{宝钗听 了，明知早晨他嘴里念诗自然是为黛玉的事了，想来他那个呆性是不能劝的，倒好 叫他睡两夜，索性自己死了心也罢了，况兼昨夜听他睡的倒也安静。}
}
{
    \eR{[English source not present in the checked-in Bencraft/Joly files.]}
}
{
}

\Verse{85}
{
    \zR{便道：“好没 来由，你只管睡去，我们拦你作什么？}
}
{
    \eR{[English source not present in the checked-in Bencraft/Joly files.]}
}
{
}

\Verse{86}
{
    \zR{但只别胡思乱想的招出些邪魔外祟来。”}
}
{
    \eR{[English source not present in the checked-in Bencraft/Joly files.]}
}
{
}

\Verse{87}
{
    \zR{宝 玉笑道：“谁想什么。”}
}
{
    \eR{[English source not present in the checked-in Bencraft/Joly files.]}
}
{
}

\Verse{88}
{
    \zR{袭人道：“依我劝，二爷竟还是屋里睡罢。}
}
{
    \eR{[English source not present in the checked-in Bencraft/Joly files.]}
}
{
}

\Verse{89}
{
    \zR{外边一时照应 不到，着了凉，倒不好。”}
}
{
    \eR{[English source not present in the checked-in Bencraft/Joly files.]}
}
{
}

\Verse{90}
{
    \zR{宝玉未及答言，宝钗却向袭人使了个眼色儿。}
}
{
    \eR{[English source not present in the checked-in Bencraft/Joly files.]}
}
{
}

\Verse{91}
{
    \zR{袭人会意， 道：“也罢，叫个人跟着你罢，夜里好倒茶倒水的。”}
}
{
    \eR{[English source not present in the checked-in Bencraft/Joly files.]}
}
{
}

\Verse{92}
{
    \zR{宝玉便笑道：“这么说，你 就跟了我来。”}
}
{
    \eR{[English source not present in the checked-in Bencraft/Joly files.]}
}
{
}

\Verse{93}
{
    \zR{袭人听了，倒没意思起来，登时飞红了脸，一声也不言语。}
}
{
    \eR{[English source not present in the checked-in Bencraft/Joly files.]}
}
{
}

\Verse{94}
{
    \zR{宝钗素 知袭人稳重，便说道：“他是跟惯了我的，还叫他跟着我罢。}
}
{
    \eR{[English source not present in the checked-in Bencraft/Joly files.]}
}
{
}

\Verse{95}
{
    \zR{叫麝月五儿照料着也 罢了。}
}
{
    \eR{[English source not present in the checked-in Bencraft/Joly files.]}
}
{
}

\Verse{96}
{
    \zR{况且今日他跟着我闹了一天，也乏了，该叫他歇歇了。”}
}
{
    \eR{[English source not present in the checked-in Bencraft/Joly files.]}
}
{
}

\Verse{97}
{
    \zR{宝玉只得笑着出来。}
}
{
    \eR{[English source not present in the checked-in Bencraft/Joly files.]}
}
{
}

\Verse{98}
{
    \zR{宝钗因命麝月五儿给宝玉仍在外间铺设了，又嘱咐两个人：“醒睡些。}
}
{
    \eR{[English source not present in the checked-in Bencraft/Joly files.]}
}
{
}

\Verse{99}
{
    \zR{要茶要水， 都留点神儿。”}
}
{
    \eR{[English source not present in the checked-in Bencraft/Joly files.]}
}
{
}

\Verse{100}
{
    \zR{两个答应着。}
}
{
    \eR{[English source not present in the checked-in Bencraft/Joly files.]}
}
{
}

\Verse{101}
{
    \zR{出来看见宝玉端然坐在床上，闭目合掌，居然像个和 尚一般，两个也不敢言语，只管瞅着他笑。}
}
{
    \eR{[English source not present in the checked-in Bencraft/Joly files.]}
}
{
}

\Verse{102}
{
    \zR{宝钗又命袭人出来照应。}
}
{
    \eR{[English source not present in the checked-in Bencraft/Joly files.]}
}
{
}

\Verse{103}
{
    \zR{袭人看见这般， 却也好笑，便轻轻的叫道：“该睡了。}
}
{
    \eR{[English source not present in the checked-in Bencraft/Joly files.]}
}
{
}

\Verse{104}
{
    \zR{怎么又打起坐来了？”}
}
{
    \eR{[English source not present in the checked-in Bencraft/Joly files.]}
}
{
}

\Verse{105}
{
    \zR{宝玉睁开眼看见袭人， 便道：“你们只管睡罢，我坐一坐就睡。”}
}
{
    \eR{[English source not present in the checked-in Bencraft/Joly files.]}
}
{
}

\Verse{106}
{
    \zR{袭人道：“因为你昨日那个光景，闹的 二奶奶一夜没睡，你再这么着成什么事？”}
}
{
    \eR{[English source not present in the checked-in Bencraft/Joly files.]}
}
{
}

\Verse{107}
{
    \zR{宝玉料着自己不睡，都不肯睡，便收拾 睡下。}
}
{
    \eR{[English source not present in the checked-in Bencraft/Joly files.]}
}
{
}

\Verse{108}
{
    \zR{袭人又嘱咐了麝月等几句，才进去关门睡了。}
}
{
    \eR{[English source not present in the checked-in Bencraft/Joly files.]}
}
{
}

\Verse{109}
{
    \zR{这里麝月五儿两个人也收拾了 被褥，伺候宝玉睡着，各自歇下。}
}
{
    \eR{[English source not present in the checked-in Bencraft/Joly files.]}
}
{
}

\Verse{110}
{
    \zR{那知宝玉要睡越睡不着，见他两个人在那里打铺，忽然想起那年袭人不在家 时，晴雯麝月两个人服事，夜间麝月出去，晴雯要唬他，因为没穿衣服着了凉，后
        来还是从这个病上死的。}
}
{
    \eR{[English source not present in the checked-in Bencraft/Joly files.]}
}
{
}

\Verse{111}
{
    \zR{想到这里，一心移在晴雯身上去了。}
}
{
    \eR{[English source not present in the checked-in Bencraft/Joly files.]}
}
{
}

\Verse{112}
{
    \zR{忽又想起凤姐说五儿 给晴雯“脱了个影儿”，因将想晴雯的心又移在五儿身上。}
}
{
    \eR{[English source not present in the checked-in Bencraft/Joly files.]}
}
{
}

\Verse{113}
{
    \zR{自己假装睡着，偷偷儿 的看那五儿，越瞧越像晴雯，不觉呆性复发。}
}
{
    \eR{[English source not present in the checked-in Bencraft/Joly files.]}
}
{
}

\Verse{114}
{
    \zR{听了听里间已无声息，知是睡了；但 不知麝月睡了没有，便故意叫了两声，却不答应。}
}
{
    \eR{[English source not present in the checked-in Bencraft/Joly files.]}
}
{
}

\Verse{115}
{
    \zR{五儿听见了宝玉叫人，便问道： “二爷要什么？”}
}
{
    \eR{[English source not present in the checked-in Bencraft/Joly files.]}
}
{
}

\Verse{116}
{
    \zR{宝玉道：“我要漱漱口。”}
}
{
    \eR{[English source not present in the checked-in Bencraft/Joly files.]}
}
{
}

\Verse{117}
{
    \zR{五儿见麝月已睡，只得起来，重新剪 了蜡花，倒了一钟茶来，一手托着漱盂。}
}
{
    \eR{[English source not present in the checked-in Bencraft/Joly files.]}
}
{
}

\Verse{118}
{
    \zR{却因赶忙起来的，身上只穿着一件桃红绫 子小袄儿，松松的挽着一个？}
}
{
    \eR{[English source not present in the checked-in Bencraft/Joly files.]}
}
{
}

\Verse{119}
{
    \zR{儿。}
}
{
    \eR{[English source not present in the checked-in Bencraft/Joly files.]}
}
{
}

\Verse{120}
{
    \zR{宝玉看时，居然晴雯复生。}
}
{
    \eR{[English source not present in the checked-in Bencraft/Joly files.]}
}
{
}

\Verse{121}
{
    \zR{忽又想起晴雯说的“早 知担了虚名，也就打个正经主意了”，不觉呆呆的呆看，也不接茶。}
}
{
    \eR{[English source not present in the checked-in Bencraft/Joly files.]}
}
{
}

\Verse{122}
{
    \zR{那五儿自从芳官去后，也无心进来了。}
}
{
    \eR{[English source not present in the checked-in Bencraft/Joly files.]}
}
{
}

\Verse{123}
{
    \zR{后来听说凤姐叫他进来伏侍宝玉，竟比 宝玉盼他进来的心还急。}
}
{
    \eR{[English source not present in the checked-in Bencraft/Joly files.]}
}
{
}

\Verse{124}
{
    \zR{不想进来以后，见宝钗袭人一般尊贵稳重，看着心里实在 敬慕；又见宝玉疯疯傻傻，不似先前的丰致；又听见王夫人为女孩子们和宝玉玩笑
        都撵了，所以把那女儿的柔情和素日的痴心，一概搁起。}
}
{
    \eR{[English source not present in the checked-in Bencraft/Joly files.]}
}
{
}

\Verse{125}
{
    \zR{怎奈这位呆爷今晚把他当 作晴雯，只管爱惜起来。}
}
{
    \eR{[English source not present in the checked-in Bencraft/Joly files.]}
}
{
}

\Verse{126}
{
    \zR{那五儿早已羞得两颊红潮，又不敢大声说话，只得轻轻的 说道：“二爷，漱口啊。”}
}
{
    \eR{[English source not present in the checked-in Bencraft/Joly files.]}
}
{
}

\Verse{127}
{
    \zR{宝玉笑着接了茶在手中，也不知道漱了没有，便笑嘻嘻 的问道：“你和晴雯姐姐好不是啊？”}
}
{
    \eR{[English source not present in the checked-in Bencraft/Joly files.]}
}
{
}

\Verse{128}
{
    \zR{五儿听了，摸不着头脑，便道：“都是姐妹， 也没有什么不好的。”}
}
{
    \eR{[English source not present in the checked-in Bencraft/Joly files.]}
}
{
}

\Verse{129}
{
    \zR{宝玉又悄悄的问道：“晴雯病重了，我看他去，不是你也去 了么？”}
}
{
    \eR{[English source not present in the checked-in Bencraft/Joly files.]}
}
{
}

\Verse{130}
{
    \zR{五儿微微笑着点头儿。}
}
{
    \eR{[English source not present in the checked-in Bencraft/Joly files.]}
}
{
}

\Verse{131}
{
    \zR{宝玉道：“你听见他说什么了没有？”}
}
{
    \eR{[English source not present in the checked-in Bencraft/Joly files.]}
}
{
}

\Verse{132}
{
    \zR{五儿摇着头 儿道：“没有。”}
}
{
    \eR{[English source not present in the checked-in Bencraft/Joly files.]}
}
{
}

\Verse{133}
{
    \zR{宝玉已经忘神，便把五儿的手一拉。}
}
{
    \eR{[English source not present in the checked-in Bencraft/Joly files.]}
}
{
}

\Verse{134}
{
    \zR{五儿急的红了脸，心里乱跳， 便悄悄说道：“二爷，有什么话只管说，别拉拉扯扯的。”}
}
{
    \eR{[English source not present in the checked-in Bencraft/Joly files.]}
}
{
}

\Verse{135}
{
    \zR{宝玉才撒了手，说道： “他和我说来着：‘早知担了个虚名，也就打正经主意了。’}
}
{
    \eR{[English source not present in the checked-in Bencraft/Joly files.]}
}
{
}

\Verse{136}
{
    \zR{你怎么没听见么？”}
}
{
    \eR{[English source not present in the checked-in Bencraft/Joly files.]}
}
{
}

\Verse{137}
{
    \zR{五儿听了，这话明明是撩拨自己的意思，又不敢怎么样，便说道：“那是他自己没 脸。}
}
{
    \eR{[English source not present in the checked-in Bencraft/Joly files.]}
}
{
}

\Verse{138}
{
    \zR{这也是我们女孩儿家说得的吗？”}
}
{
    \eR{[English source not present in the checked-in Bencraft/Joly files.]}
}
{
}

\Verse{139}
{
    \zR{宝玉着急道：“你怎么也是这么个道学先生! 我看你长的和他一模一样，我才肯和你说这个话，你怎么倒拿这些话遭塌他？”}
}
{
    \eR{[English source not present in the checked-in Bencraft/Joly files.]}
}
{
}

\Verse{140}
{
    \zR{此时五儿心中也不知宝玉是怎么个意思，便说道：“夜深了，二爷睡罢，别紧 着坐着，看凉着了。}
}
{
    \eR{[English source not present in the checked-in Bencraft/Joly files.]}
}
{
}

\Verse{141}
{
    \zR{刚才奶奶和袭人姐姐怎么嘱咐来！”}
}
{
    \eR{[English source not present in the checked-in Bencraft/Joly files.]}
}
{
}

\Verse{142}
{
    \zR{宝玉道：“我不凉。”}
}
{
    \eR{[English source not present in the checked-in Bencraft/Joly files.]}
}
{
}

\Verse{143}
{
    \zR{说 到这里，忽然想起五儿没穿着大衣裳，就怕他也像晴雯着了凉，便问道：“你为什 么不穿上衣裳就过来？”}
}
{
    \eR{[English source not present in the checked-in Bencraft/Joly files.]}
}
{
}

\Verse{144}
{
    \zR{五儿道：“爷叫的紧，那里有尽着穿衣裳的空儿？}
}
{
    \eR{[English source not present in the checked-in Bencraft/Joly files.]}
}
{
}

\Verse{145}
{
    \zR{要知道 说这半天话儿时，我也穿上了。”}
}
{
    \eR{[English source not present in the checked-in Bencraft/Joly files.]}
}
{
}

\Verse{146}
{
    \zR{宝玉听了，连忙把自己盖的一件月白绫子绵袄儿 揭起来递给五儿，叫他披上。}
}
{
    \eR{[English source not present in the checked-in Bencraft/Joly files.]}
}
{
}

\Verse{147}
{
    \zR{五儿只不肯接，说：“二爷盖着罢，我不凉。}
}
{
    \eR{[English source not present in the checked-in Bencraft/Joly files.]}
}
{
}

\Verse{148}
{
    \zR{我凉， 我有我的衣裳。”}
}
{
    \eR{[English source not present in the checked-in Bencraft/Joly files.]}
}
{
}

\Verse{149}
{
    \zR{说着，回到自己铺边，拉了一件长袄披上。}
}
{
    \eR{[English source not present in the checked-in Bencraft/Joly files.]}
}
{
}

\Verse{150}
{
    \zR{又听了听，麝月睡的 正浓，才慢慢过来说：“二爷今晚不是要养神呢吗？”}
}
{
    \eR{[English source not present in the checked-in Bencraft/Joly files.]}
}
{
}

\Verse{151}
{
    \zR{宝玉笑道：“实告诉你罢： 什么是养神!我倒是要遇仙的意思。”}
}
{
    \eR{[English source not present in the checked-in Bencraft/Joly files.]}
}
{
}

\Verse{152}
{
    \zR{五儿听了，越发动了疑心，便问道：“遇什 么仙？”}
}
{
    \eR{[English source not present in the checked-in Bencraft/Joly files.]}
}
{
}

\Verse{153}
{
    \zR{宝玉道：“你要知道，这话长着呢。}
}
{
    \eR{[English source not present in the checked-in Bencraft/Joly files.]}
}
{
}

\Verse{154}
{
    \zR{你挨着我来坐下，我告诉你。”}
}
{
    \eR{[English source not present in the checked-in Bencraft/Joly files.]}
}
{
}

\Verse{155}
{
    \zR{五儿 红了脸，笑道：“你在那里躺着，我怎么坐呢？”}
}
{
    \eR{[English source not present in the checked-in Bencraft/Joly files.]}
}
{
}

\Verse{156}
{
    \zR{宝玉道：“这个何妨？}
}
{
    \eR{[English source not present in the checked-in Bencraft/Joly files.]}
}
{
}

\Verse{157}
{
    \zR{那一年冷 天，也是你晴雯姐姐和麝月姐姐玩，我怕冻着他，还把他揽在一个被窝儿里呢。}
}
{
    \eR{[English source not present in the checked-in Bencraft/Joly files.]}
}
{
}

\Verse{158}
{
    \zR{这 有什么？}
}
{
    \eR{[English source not present in the checked-in Bencraft/Joly files.]}
}
{
}

\Verse{159}
{
    \zR{大凡一个人，总别酸文假醋的才好。”}
}
{
    \eR{[English source not present in the checked-in Bencraft/Joly files.]}
}
{
}

\Verse{160}
{
    \zR{五儿听了，句句都是宝玉调戏之意， 那知这位呆爷却是实心实意的话。}
}
{
    \eR{[English source not present in the checked-in Bencraft/Joly files.]}
}
{
}

\Verse{161}
{
    \zR{五儿此时走开不好，站着不好，坐下不好，倒没 了主意。}
}
{
    \eR{[English source not present in the checked-in Bencraft/Joly files.]}
}
{
}

\Verse{162}
{
    \zR{因拿眼一溜，抿着嘴儿笑道：“你别混说了。}
}
{
    \eR{[English source not present in the checked-in Bencraft/Joly files.]}
}
{
}

\Verse{163}
{
    \zR{看人家听见，什么意思？}
}
{
    \eR{[English source not present in the checked-in Bencraft/Joly files.]}
}
{
}

\Verse{164}
{
    \zR{怨 不得人家说你专在女孩儿身上用工夫。}
}
{
    \eR{[English source not present in the checked-in Bencraft/Joly files.]}
}
{
}

\Verse{165}
{
    \zR{你自己放着二奶奶和袭人姐姐，都是仙人儿 似的，只爱和别人混搅。}
}
{
    \eR{[English source not present in the checked-in Bencraft/Joly files.]}
}
{
}

\Verse{166}
{
    \zR{明儿再说这些话，我回了二奶奶，看你什么脸见人。”}
}
{
    \eR{[English source not present in the checked-in Bencraft/Joly files.]}
}
{
}

\Verse{167}
{
    \zR{正 说着，只听外面“咕咚”一声，把两个人吓了一跳。}
}
{
    \eR{[English source not present in the checked-in Bencraft/Joly files.]}
}
{
}

\Verse{168}
{
    \zR{里间宝钗咳嗽了一声，宝玉听 见连忙？}
}
{
    \eR{[English source not present in the checked-in Bencraft/Joly files.]}
}
{
}

\Verse{169}
{
    \zR{嘴儿，五儿也就忙忙的息了灯，悄悄的躺下了。}
}
{
    \eR{[English source not present in the checked-in Bencraft/Joly files.]}
}
{
}

\Verse{170}
{
    \zR{原来宝钗袭人因昨夜不曾 睡，又兼日间劳乏了一天，所以睡去，都不曾听见他们说话，此时院中一响，猛然 惊醒，听了听，也无动静。}
}
{
    \eR{[English source not present in the checked-in Bencraft/Joly files.]}
}
{
}

\Verse{171}
{
    \zR{宝玉此时躺在床上，心里疑惑：“莫非林妹妹来了，听 见我和五儿说话，故意吓我们的？”}
}
{
    \eR{[English source not present in the checked-in Bencraft/Joly files.]}
}
{
}

\Verse{172}
{
    \zR{翻来覆去，胡思乱想，五更以后，才朦胧睡去。}
}
{
    \eR{[English source not present in the checked-in Bencraft/Joly files.]}
}
{
}

\Verse{173}
{
    \zR{却说五儿被宝玉鬼混了半夜，又兼宝钗咳嗽，自己怀着鬼胎，生怕宝钗听见了， 也是思前想后，一夜无眠。}
}
{
    \eR{[English source not present in the checked-in Bencraft/Joly files.]}
}
{
}

\Verse{174}
{
    \zR{次日一早起来，见宝玉尚自昏昏睡着，便轻轻儿的收拾 了屋子。}
}
{
    \eR{[English source not present in the checked-in Bencraft/Joly files.]}
}
{
}

\Verse{175}
{
    \zR{那时麝月已醒，便道：“你怎么这么早起来了？}
}
{
    \eR{[English source not present in the checked-in Bencraft/Joly files.]}
}
{
}

\Verse{176}
{
    \zR{你难道一夜没睡吗？”}
}
{
    \eR{[English source not present in the checked-in Bencraft/Joly files.]}
}
{
}

\Verse{177}
{
    \zR{五 儿听这话又似麝月知道了的光景，便只是讪笑，也不答言。}
}
{
    \eR{[English source not present in the checked-in Bencraft/Joly files.]}
}
{
}

\Verse{178}
{
    \zR{一时宝钗袭人也都起来， 开了门。}
}
{
    \eR{[English source not present in the checked-in Bencraft/Joly files.]}
}
{
}

\Verse{179}
{
    \zR{见宝玉尚睡，却也纳闷：“怎么在外头两夜睡的倒这么安稳呢？”}
}
{
    \eR{[English source not present in the checked-in Bencraft/Joly files.]}
}
{
}

\Verse{180}
{
    \zR{及宝玉 醒来，见众人都起来了，自己连忙爬起。}
}
{
    \eR{[English source not present in the checked-in Bencraft/Joly files.]}
}
{
}

\Verse{181}
{
    \zR{揉着眼睛，细想昨夜又不曾梦见，可是“仙 凡路隔”了。}
}
{
    \eR{[English source not present in the checked-in Bencraft/Joly files.]}
}
{
}

\Verse{182}
{
    \zR{慢慢的下了床，又想昨夜五儿说的“宝钗袭人都是天仙一般”，这话 却也不错，便怔怔的瞅着宝钗。}
}
{
    \eR{[English source not present in the checked-in Bencraft/Joly files.]}
}
{
}

\Verse{183}
{
    \zR{宝钗见他发怔，虽知他为黛玉之事，却也定不得梦不梦，只是瞅的自己倒不好 意思的，便道：“你昨夜可遇见仙了么？”}
}
{
    \eR{[English source not present in the checked-in Bencraft/Joly files.]}
}
{
}

\Verse{184}
{
    \zR{宝玉听了，只道昨晚的话宝钗听见了， 笑着勉强说道：“这是那里的话？”}
}
{
    \eR{[English source not present in the checked-in Bencraft/Joly files.]}
}
{
}

\Verse{185}
{
    \zR{那五儿听了这一句，越发心虚起来，又不好说 的，只得且看宝钗的光景。}
}
{
    \eR{[English source not present in the checked-in Bencraft/Joly files.]}
}
{
}

\Verse{186}
{
    \zR{只见宝钗又笑着问五儿道：“你听见二爷睡梦里和人说 话来着么？”}
}
{
    \eR{[English source not present in the checked-in Bencraft/Joly files.]}
}
{
}

\Verse{187}
{
    \zR{宝玉听了，自己坐不住，搭讪着走开了。}
}
{
    \eR{[English source not present in the checked-in Bencraft/Joly files.]}
}
{
}

\Verse{188}
{
    \zR{五儿把脸飞红，只得含糊道： “前半夜倒说了几句，我也没听真。}
}
{
    \eR{[English source not present in the checked-in Bencraft/Joly files.]}
}
{
}

\Verse{189}
{
    \zR{什么‘担了虚名’，又什么‘没打正经主意’， 我也不懂，劝着二爷睡了。}
}
{
    \eR{[English source not present in the checked-in Bencraft/Joly files.]}
}
{
}

\Verse{190}
{
    \zR{后来我也睡了，不知二爷还说来着没有。”}
}
{
    \eR{[English source not present in the checked-in Bencraft/Joly files.]}
}
{
}

\Verse{191}
{
    \zR{宝钗低头一 想：“这话明是为黛玉了。}
}
{
    \eR{[English source not present in the checked-in Bencraft/Joly files.]}
}
{
}

\Verse{192}
{
    \zR{但尽着叫他在外头，恐怕心邪了，招出些花妖柳怪来。}
}
{
    \eR{[English source not present in the checked-in Bencraft/Joly files.]}
}
{
}

\Verse{193}
{
    \zR{况兼他的旧病，原在姐妹上情重，只好设法将他的心意挪移过来，然后能免无事。”}
}
{
    \eR{[English source not present in the checked-in Bencraft/Joly files.]}
}
{
}

\Verse{194}
{
    \zR{想到这里，不免面红耳热起来，也就讪讪的进房梳洗去了。}
}
{
    \eR{[English source not present in the checked-in Bencraft/Joly files.]}
}
{
}

\Verse{195}
{
    \zR{且说贾母两日高兴，略吃多了些，这晚有些不受用；第二天，便觉着胸口饱闷。}
}
{
    \eR{[English source not present in the checked-in Bencraft/Joly files.]}
}
{
}

\Verse{196}
{
    \zR{鸳鸯等要回贾政，贾母不叫言语，说：“我这两日嘴馋些，吃多了点子。}
}
{
    \eR{[English source not present in the checked-in Bencraft/Joly files.]}
}
{
}

\Verse{197}
{
    \zR{我饿一顿 就好了，你们快别吵嚷。”}
}
{
    \eR{[English source not present in the checked-in Bencraft/Joly files.]}
}
{
}

\Verse{198}
{
    \zR{于是鸳鸯等并没有告诉人。}
}
{
    \eR{[English source not present in the checked-in Bencraft/Joly files.]}
}
{
}

\Verse{199}
{
    \zR{这日晚间，宝玉回到自己屋 里，见宝钗自贾母王夫人处才请了晚安回来。}
}
{
    \eR{[English source not present in the checked-in Bencraft/Joly files.]}
}
{
}

\Verse{200}
{
    \zR{宝玉想着早起之事，未免赧颜抱惭， 宝钗看他这样的，也晓得是没意思的光景。}
}
{
    \eR{[English source not present in the checked-in Bencraft/Joly files.]}
}
{
}

\Verse{201}
{
    \zR{因想着他是个痴情人，要治他的这个病， 少不得仍以痴情治之。}
}
{
    \eR{[English source not present in the checked-in Bencraft/Joly files.]}
}
{
}

\Verse{202}
{
    \zR{想了想，便问宝玉道：“你今夜还在外头睡去罢咧？”}
}
{
    \eR{[English source not present in the checked-in Bencraft/Joly files.]}
}
{
}

\Verse{203}
{
    \zR{宝玉 自觉没趣，便道：“里头外头都是一样的。”}
}
{
    \eR{[English source not present in the checked-in Bencraft/Joly files.]}
}
{
}

\Verse{204}
{
    \zR{宝钗意欲再说，反觉碍难出口。}
}
{
    \eR{[English source not present in the checked-in Bencraft/Joly files.]}
}
{
}

\Verse{205}
{
    \zR{袭人 道：“罢呀，这倒是什么道理呢？}
}
{
    \eR{[English source not present in the checked-in Bencraft/Joly files.]}
}
{
}

\Verse{206}
{
    \zR{我不信睡的那么安顿。”}
}
{
    \eR{[English source not present in the checked-in Bencraft/Joly files.]}
}
{
}

\Verse{207}
{
    \zR{五儿听见这话，连忙接 口道：“二爷在外头睡，别的倒没有什么，只爱说梦话，叫人摸不着头脑儿，又不 敢驳他的回。”}
}
{
    \eR{[English source not present in the checked-in Bencraft/Joly files.]}
}
{
}

\Verse{208}
{
    \zR{袭人便道：“我今日挪出床上睡睡，看说梦话不说。}
}
{
    \eR{[English source not present in the checked-in Bencraft/Joly files.]}
}
{
}

\Verse{209}
{
    \zR{你们只管把二 爷的铺盖铺在里间就完了。”}
}
{
    \eR{[English source not present in the checked-in Bencraft/Joly files.]}
}
{
}

\Verse{210}
{
    \zR{宝钗听了，也不作声。}
}
{
    \eR{[English source not present in the checked-in Bencraft/Joly files.]}
}
{
}

\Verse{211}
{
    \zR{宝玉自己惭愧，那里还有强嘴 的分儿，便依着搬进来。}
}
{
    \eR{[English source not present in the checked-in Bencraft/Joly files.]}
}
{
}

\Verse{212}
{
    \zR{一则宝玉抱歉，欲安宝钗之心；二则宝钗恐宝玉思郁成疾， 不如稍示柔情，使得亲近，以为移花接木之计。}
}
{
    \eR{[English source not present in the checked-in Bencraft/Joly files.]}
}
{
}

\Verse{213}
{
    \zR{于是当晚袭人果然挪出去。}
}
{
    \eR{[English source not present in the checked-in Bencraft/Joly files.]}
}
{
}

\Verse{214}
{
    \zR{这宝玉 固然是有意负荆，那宝钗自然也无心拒客，从过门至今日，方才是雨腻云香，氤氲 调畅。}
}
{
    \eR{[English source not present in the checked-in Bencraft/Joly files.]}
}
{
}

\Verse{215}
{
    \zR{从此“二五之精，妙合而凝”。}
}
{
    \eR{[English source not present in the checked-in Bencraft/Joly files.]}
}
{
}

\Verse{216}
{
    \zR{此是后话不提。}
}
{
    \eR{[English source not present in the checked-in Bencraft/Joly files.]}
}
{
}

\Verse{217}
{
    \zR{且说次日宝玉宝钗同起，宝玉梳洗了，先过贾母这边来。}
}
{
    \eR{[English source not present in the checked-in Bencraft/Joly files.]}
}
{
}

\Verse{218}
{
    \zR{这里贾母因疼宝玉， 又想宝钗孝顺，忽然想起一件东西来。}
}
{
    \eR{[English source not present in the checked-in Bencraft/Joly files.]}
}
{
}

\Verse{219}
{
    \zR{便叫鸳鸯开了箱子，取出祖上所遗的一个汉 玉？}
}
{
    \eR{[English source not present in the checked-in Bencraft/Joly files.]}
}
{
}

\Verse{220}
{
    \zR{，虽不及宝玉他那块玉石，挂在身上却也希罕。}
}
{
    \eR{[English source not present in the checked-in Bencraft/Joly files.]}
}
{
}

\Verse{221}
{
    \zR{鸳鸯找出来递与贾母，便说道： “这件东西，我好像从没见的。}
}
{
    \eR{[English source not present in the checked-in Bencraft/Joly files.]}
}
{
}

\Verse{222}
{
    \zR{老太太这些年还记得这样清楚，说是那一箱什么匣 子里装着，我按着老太太的话一拿就拿出来了。}
}
{
    \eR{[English source not present in the checked-in Bencraft/Joly files.]}
}
{
}

\Verse{223}
{
    \zR{老太太这会子叫拿出来做什么？”}
}
{
    \eR{[English source not present in the checked-in Bencraft/Joly files.]}
}
{
}

\Verse{224}
{
    \zR{贾母道：“你那里知道？}
}
{
    \eR{[English source not present in the checked-in Bencraft/Joly files.]}
}
{
}

\Verse{225}
{
    \zR{这块玉还是祖爷爷给我们老太爷，老太爷疼我，临出嫁的 时候叫了我去，亲手递给我的。}
}
{
    \eR{[English source not present in the checked-in Bencraft/Joly files.]}
}
{
}

\Verse{226}
{
    \zR{还说：‘这玉是汉朝所佩的东西，很贵重，你拿着 就像见了我的一样。’}
}
{
    \eR{[English source not present in the checked-in Bencraft/Joly files.]}
}
{
}

\Verse{227}
{
    \zR{我那时还小，拿了来也不当什么便撩在箱子里。}
}
{
    \eR{[English source not present in the checked-in Bencraft/Joly files.]}
}
{
}

\Verse{228}
{
    \zR{到了这里， 我见咱们家的东西也多，这算得什么，从没带过，一撩便撩了六十多年。}
}
{
    \eR{[English source not present in the checked-in Bencraft/Joly files.]}
}
{
}

\Verse{229}
{
    \zR{今儿见宝 玉这样孝顺，他又丢了一块玉，故此想着拿出来给他，也像是祖上给我的意思。”}
}
{
    \eR{[English source not present in the checked-in Bencraft/Joly files.]}
}
{
}

\Verse{230}
{
    \zR{一时宝玉请了安，贾母便喜欢道：“你过来，我给你一件东西瞧瞧。”}
}
{
    \eR{[English source not present in the checked-in Bencraft/Joly files.]}
}
{
}

\Verse{231}
{
    \zR{宝玉走到床 前，贾母便把那块汉玉递给宝玉。}
}
{
    \eR{[English source not present in the checked-in Bencraft/Joly files.]}
}
{
}

\Verse{232}
{
    \zR{宝玉接来一瞧，那玉有三寸方圆，形似甜瓜，色 有红晕，甚是精致。}
}
{
    \eR{[English source not present in the checked-in Bencraft/Joly files.]}
}
{
}

\Verse{233}
{
    \zR{宝玉口口称赞。}
}
{
    \eR{[English source not present in the checked-in Bencraft/Joly files.]}
}
{
}

\Verse{234}
{
    \zR{贾母道：“你爱么？}
}
{
    \eR{[English source not present in the checked-in Bencraft/Joly files.]}
}
{
}

\Verse{235}
{
    \zR{这是我祖爷爷给我的，我 传了你罢。”}
}
{
    \eR{[English source not present in the checked-in Bencraft/Joly files.]}
}
{
}

\Verse{236}
{
    \zR{宝玉笑着，请了个安谢了，又拿了要送给他母亲瞧。}
}
{
    \eR{[English source not present in the checked-in Bencraft/Joly files.]}
}
{
}

\Verse{237}
{
    \zR{贾母道：“你太 太瞧了，告诉你老子，又说疼儿子不如疼孙子了。}
}
{
    \eR{[English source not present in the checked-in Bencraft/Joly files.]}
}
{
}

\Verse{238}
{
    \zR{他们从没见过。”}
}
{
    \eR{[English source not present in the checked-in Bencraft/Joly files.]}
}
{
}

\Verse{239}
{
    \zR{宝玉笑着去了。}
}
{
    \eR{[English source not present in the checked-in Bencraft/Joly files.]}
}
{
}

\Verse{240}
{
    \zR{宝钗等又说了几句话，也辞了出来。}
}
{
    \eR{[English source not present in the checked-in Bencraft/Joly files.]}
}
{
}

\Verse{241}
{
    \zR{自此，贾母两日不进饮食，胸口仍是膨闷，觉得头晕目眩，咳嗽。}
}
{
    \eR{[English source not present in the checked-in Bencraft/Joly files.]}
}
{
}

\Verse{242}
{
    \zR{邢王二夫人、 凤姐等请安，见贾母精神尚好，不过叫人告诉贾政，立刻来请了安。}
}
{
    \eR{[English source not present in the checked-in Bencraft/Joly files.]}
}
{
}

\Verse{243}
{
    \zR{贾政出来，即 请大夫看脉。}
}
{
    \eR{[English source not present in the checked-in Bencraft/Joly files.]}
}
{
}

\Verse{244}
{
    \zR{不多一时，大夫来诊了脉，说是有年纪的人，停了些饮食，感冒些风 寒，略消导发散些就好了。}
}
{
    \eR{[English source not present in the checked-in Bencraft/Joly files.]}
}
{
}

\Verse{245}
{
    \zR{开了方子，贾政看了，知是寻常药品，命人煎好进服。}
}
{
    \eR{[English source not present in the checked-in Bencraft/Joly files.]}
}
{
}

\Verse{246}
{
    \zR{以后贾政早晚进来请安。}
}
{
    \eR{[English source not present in the checked-in Bencraft/Joly files.]}
}
{
}

\Verse{247}
{
    \zR{一连三日，不见稍减。}
}
{
    \eR{[English source not present in the checked-in Bencraft/Joly files.]}
}
{
}

\Verse{248}
{
    \zR{贾政又命贾琏打听好大夫，“快去 请来瞧老太太的病。}
}
{
    \eR{[English source not present in the checked-in Bencraft/Joly files.]}
}
{
}

\Verse{249}
{
    \zR{咱们家常请的几个大夫，我瞧着不怎么好，所以叫你去。”}
}
{
    \eR{[English source not present in the checked-in Bencraft/Joly files.]}
}
{
}

\Verse{250}
{
    \zR{贾 琏想了一想，说道：“记得那年宝兄弟病的时候，倒是请了一个不行医的来瞧好了 的，如今不如找他。”}
}
{
    \eR{[English source not present in the checked-in Bencraft/Joly files.]}
}
{
}

\Verse{251}
{
    \zR{贾政道：“医道却是极难的，越是不兴时的大夫倒有本领。}
}
{
    \eR{[English source not present in the checked-in Bencraft/Joly files.]}
}
{
}

\Verse{252}
{
    \zR{你就打发人去找来罢。”}
}
{
    \eR{[English source not present in the checked-in Bencraft/Joly files.]}
}
{
}

\Verse{253}
{
    \zR{贾琏即忙答应去了，回来说道：“这刘大夫新近出城教书 去了，过十来天进城一次。}
}
{
    \eR{[English source not present in the checked-in Bencraft/Joly files.]}
}
{
}

\Verse{254}
{
    \zR{这时等不得，又请了一位，也就来了。”}
}
{
    \eR{[English source not present in the checked-in Bencraft/Joly files.]}
}
{
}

\Verse{255}
{
    \zR{贾政听了，只 得等着，不提。}
}
{
    \eR{[English source not present in the checked-in Bencraft/Joly files.]}
}
{
}

\Verse{256}
{
    \zR{且说贾母病时，合宅女眷无日不来请安。}
}
{
    \eR{[English source not present in the checked-in Bencraft/Joly files.]}
}
{
}

\Verse{257}
{
    \zR{一日，众人都在那里，只见看园内腰 门的老婆子进来回说：“园里的栊翠庵的妙师父知道老太太病了，特来请安。”}
}
{
    \eR{[English source not present in the checked-in Bencraft/Joly files.]}
}
{
}

\Verse{258}
{
    \zR{众 人道：“他不常过来，今儿特来，你们快请进来。”}
}
{
    \eR{[English source not present in the checked-in Bencraft/Joly files.]}
}
{
}

\Verse{259}
{
    \zR{凤姐走到床前回了贾母。}
}
{
    \eR{[English source not present in the checked-in Bencraft/Joly files.]}
}
{
}

\Verse{260}
{
    \zR{岫烟 是妙玉的旧相识，先走出去接他。}
}
{
    \eR{[English source not present in the checked-in Bencraft/Joly files.]}
}
{
}

\Verse{261}
{
    \zR{只见妙玉头带妙常冠，身上穿一件月白素绸袄儿， 外罩一件水田青缎镶边长背心，拴着秋香色的丝绦，腰下系一条淡墨画的白绫裙，
        手执麈尾念珠，跟着一个侍儿，飘飘拽拽的走来。}
}
{
    \eR{[English source not present in the checked-in Bencraft/Joly files.]}
}
{
}

\Verse{262}
{
    \zR{岫烟见了问好，说是：“在园内 住的时候儿，可以常来瞧瞧你；近来因为园内人少，一个人轻易难出来。}
}
{
    \eR{[English source not present in the checked-in Bencraft/Joly files.]}
}
{
}

\Verse{263}
{
    \zR{况且咱们 这里的腰门常关着，所以这些日子不得见你。}
}
{
    \eR{[English source not present in the checked-in Bencraft/Joly files.]}
}
{
}

\Verse{264}
{
    \zR{今儿幸会。”}
}
{
    \eR{[English source not present in the checked-in Bencraft/Joly files.]}
}
{
}

\Verse{265}
{
    \zR{妙玉道：“头里你们是 热闹场中，你们虽在外园里住，我也不便常来亲近。}
}
{
    \eR{[English source not present in the checked-in Bencraft/Joly files.]}
}
{
}

\Verse{266}
{
    \zR{如今知道这里的事情也不大好， 又听说是老太太病着，又惦记着你，还要瞧瞧宝姑娘。}
}
{
    \eR{[English source not present in the checked-in Bencraft/Joly files.]}
}
{
}

\Verse{267}
{
    \zR{我那管你们关不关？}
}
{
    \eR{[English source not present in the checked-in Bencraft/Joly files.]}
}
{
}

\Verse{268}
{
    \zR{我要来 就来，我不来，你们要我来也不能啊。”}
}
{
    \eR{[English source not present in the checked-in Bencraft/Joly files.]}
}
{
}

\Verse{269}
{
    \zR{岫烟笑道：“你还是这种脾气。”}
}
{
    \eR{[English source not present in the checked-in Bencraft/Joly files.]}
}
{
}

\Verse{270}
{
    \zR{一面说着，已到贾母房中。}
}
{
    \eR{[English source not present in the checked-in Bencraft/Joly files.]}
}
{
}

\Verse{271}
{
    \zR{众人见了，都问了好。}
}
{
    \eR{[English source not present in the checked-in Bencraft/Joly files.]}
}
{
}

\Verse{272}
{
    \zR{妙玉走到贾母床前问候，说 了几句套话。}
}
{
    \eR{[English source not present in the checked-in Bencraft/Joly files.]}
}
{
}

\Verse{273}
{
    \zR{贾母便道：“你是个女菩萨，你瞧瞧我的病可好的了好不了？”}
}
{
    \eR{[English source not present in the checked-in Bencraft/Joly files.]}
}
{
}

\Verse{274}
{
    \zR{妙玉 道：“老太太这样慈善的人，寿数正有呢。}
}
{
    \eR{[English source not present in the checked-in Bencraft/Joly files.]}
}
{
}

\Verse{275}
{
    \zR{一时感冒，吃几帖药，想来也就好了。}
}
{
    \eR{[English source not present in the checked-in Bencraft/Joly files.]}
}
{
}

\Verse{276}
{
    \zR{有年纪的人，只要宽心些。”}
}
{
    \eR{[English source not present in the checked-in Bencraft/Joly files.]}
}
{
}

\Verse{277}
{
    \zR{贾母道：“我倒不为这些。}
}
{
    \eR{[English source not present in the checked-in Bencraft/Joly files.]}
}
{
}

\Verse{278}
{
    \zR{我是极爱寻快乐的。}
}
{
    \eR{[English source not present in the checked-in Bencraft/Joly files.]}
}
{
}

\Verse{279}
{
    \zR{如今 这病也不觉怎么着，只是胸膈饱闷。}
}
{
    \eR{[English source not present in the checked-in Bencraft/Joly files.]}
}
{
}

\Verse{280}
{
    \zR{刚才大夫说是气恼所致。}
}
{
    \eR{[English source not present in the checked-in Bencraft/Joly files.]}
}
{
}

\Verse{281}
{
    \zR{你是知道的，谁敢给 我气受？}
}
{
    \eR{[English source not present in the checked-in Bencraft/Joly files.]}
}
{
}

\Verse{282}
{
    \zR{这不是那大夫脉理平常么？}
}
{
    \eR{[English source not present in the checked-in Bencraft/Joly files.]}
}
{
}

\Verse{283}
{
    \zR{我和琏儿说了，还是头一个大夫说感冒伤食的 是，明儿还请他来。”}
}
{
    \eR{[English source not present in the checked-in Bencraft/Joly files.]}
}
{
}

\Verse{284}
{
    \zR{说着，叫鸳鸯：“吩咐厨房里办一桌净素菜来，请妙师父这 里便饭。”}
}
{
    \eR{[English source not present in the checked-in Bencraft/Joly files.]}
}
{
}

\Verse{285}
{
    \zR{妙玉道：“我吃过午饭了，我是不吃东西的。”}
}
{
    \eR{[English source not present in the checked-in Bencraft/Joly files.]}
}
{
}

\Verse{286}
{
    \zR{王夫人道：“不吃也罢， 咱们多坐一会，说些闲话儿罢。”}
}
{
    \eR{[English source not present in the checked-in Bencraft/Joly files.]}
}
{
}

\Verse{287}
{
    \zR{妙玉道：“我久已不见你们，今日来瞧瞧。”}
}
{
    \eR{[English source not present in the checked-in Bencraft/Joly files.]}
}
{
}

\Verse{288}
{
    \zR{又 说了一回话，便要走。}
}
{
    \eR{[English source not present in the checked-in Bencraft/Joly files.]}
}
{
}

\Verse{289}
{
    \zR{回头见惜春站着，便问道：“四姑娘为什么这样瘦？}
}
{
    \eR{[English source not present in the checked-in Bencraft/Joly files.]}
}
{
}

\Verse{290}
{
    \zR{不要只 管爱画劳了心。”}
}
{
    \eR{[English source not present in the checked-in Bencraft/Joly files.]}
}
{
}

\Verse{291}
{
    \zR{惜春道：“我久不画了。}
}
{
    \eR{[English source not present in the checked-in Bencraft/Joly files.]}
}
{
}

\Verse{292}
{
    \zR{如今住的房屋不比园里的显亮，所以没 兴头画。”}
}
{
    \eR{[English source not present in the checked-in Bencraft/Joly files.]}
}
{
}

\Verse{293}
{
    \zR{妙玉道：“你如今住在那一所？”}
}
{
    \eR{[English source not present in the checked-in Bencraft/Joly files.]}
}
{
}

\Verse{294}
{
    \zR{惜春道：“就是你才来的那个门东边 的屋子，你要来很近。”}
}
{
    \eR{[English source not present in the checked-in Bencraft/Joly files.]}
}
{
}

\Verse{295}
{
    \zR{妙玉道：“我高兴的时候来瞧你。”}
}
{
    \eR{[English source not present in the checked-in Bencraft/Joly files.]}
}
{
}

\Verse{296}
{
    \zR{惜春等说着送了出去。}
}
{
    \eR{[English source not present in the checked-in Bencraft/Joly files.]}
}
{
}

\Verse{297}
{
    \zR{回身过来，听见丫头们回说大夫在贾母那边呢，众人暂且散去。}
}
{
    \eR{[English source not present in the checked-in Bencraft/Joly files.]}
}
{
}

\Verse{298}
{
    \zR{那知贾母这病日重一日，延医调治不效，以后又添腹泻。}
}
{
    \eR{[English source not present in the checked-in Bencraft/Joly files.]}
}
{
}

\Verse{299}
{
    \zR{贾政着急，知病难医， 即命人到衙门告诉，日夜同王夫人亲侍汤药。}
}
{
    \eR{[English source not present in the checked-in Bencraft/Joly files.]}
}
{
}

\Verse{300}
{
    \zR{一日，见贾母略进些饮食，心里稍宽， 只见老婆子在门外探头。}
}
{
    \eR{[English source not present in the checked-in Bencraft/Joly files.]}
}
{
}

\Verse{301}
{
    \zR{王夫人叫彩云看去，问问是谁。}
}
{
    \eR{[English source not present in the checked-in Bencraft/Joly files.]}
}
{
}

\Verse{302}
{
    \zR{彩云看了是陪迎春到孙家 去的人，便道：“你来做什么？”}
}
{
    \eR{[English source not present in the checked-in Bencraft/Joly files.]}
}
{
}

\Verse{303}
{
    \zR{婆子道：“我来了半日，这里找不着一个姐姐们， 我又不敢冒撞，我心里又急。”}
}
{
    \eR{[English source not present in the checked-in Bencraft/Joly files.]}
}
{
}

\Verse{304}
{
    \zR{彩云道：“你急什么？}
}
{
    \eR{[English source not present in the checked-in Bencraft/Joly files.]}
}
{
}

\Verse{305}
{
    \zR{又是姑爷作践姑娘不成么？”}
}
{
    \eR{[English source not present in the checked-in Bencraft/Joly files.]}
}
{
}

\Verse{306}
{
    \zR{婆子道：“姑娘不好了!前儿闹了一场，姑娘哭了一夜，昨日痰堵住了。}
}
{
    \eR{[English source not present in the checked-in Bencraft/Joly files.]}
}
{
}

\Verse{307}
{
    \zR{他们又不 请大夫，今日更利害了。”}
}
{
    \eR{[English source not present in the checked-in Bencraft/Joly files.]}
}
{
}

\Verse{308}
{
    \zR{彩云道：“老太太病着呢，别大惊小怪的。”}
}
{
    \eR{[English source not present in the checked-in Bencraft/Joly files.]}
}
{
}

\Verse{309}
{
    \zR{王夫人在 内已听见了，恐老太太听见不受用，忙叫彩云带他外头说去。}
}
{
    \eR{[English source not present in the checked-in Bencraft/Joly files.]}
}
{
}

\Verse{310}
{
    \zR{岂知贾母病中心静， 偏偏听见，便道：“迎丫头要死了么？”}
}
{
    \eR{[English source not present in the checked-in Bencraft/Joly files.]}
}
{
}

\Verse{311}
{
    \zR{王夫人便道：“没有。}
}
{
    \eR{[English source not present in the checked-in Bencraft/Joly files.]}
}
{
}

\Verse{312}
{
    \zR{婆子们不知轻重， 说是这两日有些病，恐不能就好，到这里问大夫。”}
}
{
    \eR{[English source not present in the checked-in Bencraft/Joly files.]}
}
{
}

\Verse{313}
{
    \zR{贾母道：“瞧我的大夫就好， 快请了去。”}
}
{
    \eR{[English source not present in the checked-in Bencraft/Joly files.]}
}
{
}

\Verse{314}
{
    \zR{王夫人便叫彩云：“叫这婆子去回大太太去。”}
}
{
    \eR{[English source not present in the checked-in Bencraft/Joly files.]}
}
{
}

\Verse{315}
{
    \zR{那婆子去了。}
}
{
    \eR{[English source not present in the checked-in Bencraft/Joly files.]}
}
{
}

\Verse{316}
{
    \zR{这里贾 母便悲伤起来，说是：“我三个孙女儿：一个享尽了福死了；三丫头远嫁，不得见
        面；迎丫头虽苦，或者熬出来，不打量他年轻轻儿的就要死了!留着我这么大年纪 的人活着做什么！”}
}
{
    \eR{[English source not present in the checked-in Bencraft/Joly files.]}
}
{
}

\Verse{317}
{
    \zR{王夫人鸳鸯等解劝了好半天。}
}
{
    \eR{[English source not present in the checked-in Bencraft/Joly files.]}
}
{
}

\Verse{318}
{
    \zR{那时宝钗李氏等不在房中，凤姐 近来有病，王夫人恐贾母生悲添病，便叫人叫了他们来陪着，自己回到房中，叫彩
        云来埋怨：“这婆子不懂事!以后我在老太太那里，你们有事，不用来回。”}
}
{
    \eR{[English source not present in the checked-in Bencraft/Joly files.]}
}
{
}

\Verse{319}
{
    \zR{丫头 们依命不言。}
}
{
    \eR{[English source not present in the checked-in Bencraft/Joly files.]}
}
{
}

\Verse{320}
{
    \zR{岂知那婆子刚到邢夫人那里，外头的人已传进来，说：“二姑奶奶死 了。”}
}
{
    \eR{[English source not present in the checked-in Bencraft/Joly files.]}
}
{
}

\Verse{321}
{
    \zR{邢夫人听了，也便哭了一场。}
}
{
    \eR{[English source not present in the checked-in Bencraft/Joly files.]}
}
{
}

\Verse{322}
{
    \zR{现今他父亲不在家中，只得叫贾琏快去瞧看。}
}
{
    \eR{[English source not present in the checked-in Bencraft/Joly files.]}
}
{
}

\Verse{323}
{
    \zR{知贾母病重，众人都不敢回。}
}
{
    \eR{[English source not present in the checked-in Bencraft/Joly files.]}
}
{
}

\Verse{324}
{
    \zR{可怜一位如花似月之女，结缡年馀，不料被孙家揉搓， 以致身亡。}
}
{
    \eR{[English source not present in the checked-in Bencraft/Joly files.]}
}
{
}

\Verse{325}
{
    \zR{又值贾母病笃，众人不便离开，竟容孙家草草完结。}
}
{
    \eR{[English source not present in the checked-in Bencraft/Joly files.]}
}
{
}

\Verse{326}
{
    \zR{贾母病势日增，只想这些孙女儿。}
}
{
    \eR{[English source not present in the checked-in Bencraft/Joly files.]}
}
{
}

\Verse{327}
{
    \zR{一时想起湘云，便打发人去瞧他。}
}
{
    \eR{[English source not present in the checked-in Bencraft/Joly files.]}
}
{
}

\Verse{328}
{
    \zR{回来的人 悄悄的找鸳鸯。}
}
{
    \eR{[English source not present in the checked-in Bencraft/Joly files.]}
}
{
}

\Verse{329}
{
    \zR{因鸳鸯在老太太身旁，王夫人等都在那里，不便上去，到了后头， 找了琥珀，告诉他道：“老太太想史姑娘，叫我们去打听。}
}
{
    \eR{[English source not present in the checked-in Bencraft/Joly files.]}
}
{
}

\Verse{330}
{
    \zR{那里知道史姑娘哭的了 不得，说是姑爷得了暴病，大夫都瞧了，说这病只怕不能好，若是变了痨病，还可 捱个四五年。}
}
{
    \eR{[English source not present in the checked-in Bencraft/Joly files.]}
}
{
}

\Verse{331}
{
    \zR{所以史姑娘心里着急。}
}
{
    \eR{[English source not present in the checked-in Bencraft/Joly files.]}
}
{
}

\Verse{332}
{
    \zR{又知道老太太病，只是不能过来请安。}
}
{
    \eR{[English source not present in the checked-in Bencraft/Joly files.]}
}
{
}

\Verse{333}
{
    \zR{还叫我 别在老太太跟前提起来，倘或老太太问起来，务必托你们变个法儿回老太太才好。”}
}
{
    \eR{[English source not present in the checked-in Bencraft/Joly files.]}
}
{
}

\Verse{334}
{
    \zR{琥珀听了，？}
}
{
    \eR{[English source not present in the checked-in Bencraft/Joly files.]}
}
{
}

\Verse{335}
{
    \zR{了一声，也就不言语了，半日说道：“你去罢。”}
}
{
    \eR{[English source not present in the checked-in Bencraft/Joly files.]}
}
{
}

\Verse{336}
{
    \zR{琥珀也不便回，心 里打算告诉鸳鸯叫他撒谎去，所以来到贾母床前。}
}
{
    \eR{[English source not present in the checked-in Bencraft/Joly files.]}
}
{
}

\Verse{337}
{
    \zR{见贾母神色大变，地下站着一屋 子的人，嘁嘁喳喳的说：“瞧着是不好。”}
}
{
    \eR{[English source not present in the checked-in Bencraft/Joly files.]}
}
{
}

\Verse{338}
{
    \zR{也不敢言语了。}
}
{
    \eR{[English source not present in the checked-in Bencraft/Joly files.]}
}
{
}

\Verse{339}
{
    \zR{这里贾政悄悄的叫贾琏 到身旁，向耳边说了几句话。}
}
{
    \eR{[English source not present in the checked-in Bencraft/Joly files.]}
}
{
}

\Verse{340}
{
    \zR{贾琏轻轻的答应，出去了，便传齐了现在家里的一干 人，说：“老太太的事，待好出来了，你们快快分头派人办去。}
}
{
    \eR{[English source not present in the checked-in Bencraft/Joly files.]}
}
{
}

\Verse{341}
{
    \zR{头一件，先请出板 来瞧瞧，好挂里子。}
}
{
    \eR{[English source not present in the checked-in Bencraft/Joly files.]}
}
{
}

\Verse{342}
{
    \zR{快到各处将各人的衣服量了尺寸，都开明了，便叫裁缝去做孝 衣。}
}
{
    \eR{[English source not present in the checked-in Bencraft/Joly files.]}
}
{
}

\Verse{343}
{
    \zR{那棚杠执事都讲定了。}
}
{
    \eR{[English source not present in the checked-in Bencraft/Joly files.]}
}
{
}

\Verse{344}
{
    \zR{厨房里还该多派几个人。”}
}
{
    \eR{[English source not present in the checked-in Bencraft/Joly files.]}
}
{
}

\Verse{345}
{
    \zR{赖大等回道：“二爷，这些 事不用爷费心，我们早打算好了，只是这项银子在那里领呢？”}
}
{
    \eR{[English source not present in the checked-in Bencraft/Joly files.]}
}
{
}

\Verse{346}
{
    \zR{贾琏道：“这种银 子不用外头去，老太太自己早留下了。}
}
{
    \eR{[English source not present in the checked-in Bencraft/Joly files.]}
}
{
}

\Verse{347}
{
    \zR{刚才老爷的主意，只要办的好，我想外面也 要好看。”}
}
{
    \eR{[English source not present in the checked-in Bencraft/Joly files.]}
}
{
}

\Verse{348}
{
    \zR{赖大等答应，派人分头办去。}
}
{
    \eR{[English source not present in the checked-in Bencraft/Joly files.]}
}
{
}

\Verse{349}
{
    \zR{贾琏复回到自己房中，便问平儿：“你奶奶今儿怎么样？”}
}
{
    \eR{[English source not present in the checked-in Bencraft/Joly files.]}
}
{
}

\Verse{350}
{
    \zR{平儿把嘴往里一努， 说：“你瞧去。”}
}
{
    \eR{[English source not present in the checked-in Bencraft/Joly files.]}
}
{
}

\Verse{351}
{
    \zR{贾琏进内，见凤姐正要穿衣，一时动不得，暂且靠在炕桌儿上。}
}
{
    \eR{[English source not present in the checked-in Bencraft/Joly files.]}
}
{
}

\Verse{352}
{
    \zR{贾琏道：“你只怕养不住了，老太太的事，今儿明儿就要出来了，你还脱得过么？}
}
{
    \eR{[English source not present in the checked-in Bencraft/Joly files.]}
}
{
}

\Verse{353}
{
    \zR{快叫人将屋里收拾收拾，就该扎挣上去了。}
}
{
    \eR{[English source not present in the checked-in Bencraft/Joly files.]}
}
{
}

\Verse{354}
{
    \zR{若有了事，你我还能回来么？”}
}
{
    \eR{[English source not present in the checked-in Bencraft/Joly files.]}
}
{
}

\Verse{355}
{
    \zR{凤姐道： “咱们这里还有什么收拾的!不过就是这点子东西，还怕什么？}
}
{
    \eR{[English source not present in the checked-in Bencraft/Joly files.]}
}
{
}

\Verse{356}
{
    \zR{你先去罢，看老爷叫 你。}
}
{
    \eR{[English source not present in the checked-in Bencraft/Joly files.]}
}
{
}

\Verse{357}
{
    \zR{我换件衣裳就来。”}
}
{
    \eR{[English source not present in the checked-in Bencraft/Joly files.]}
}
{
}

\Verse{358}
{
    \zR{贾琏先回到贾母房里，向贾政悄悄的回道：“诸事已交派 明白了。”}
}
{
    \eR{[English source not present in the checked-in Bencraft/Joly files.]}
}
{
}

\Verse{359}
{
    \zR{贾政点头。}
}
{
    \eR{[English source not present in the checked-in Bencraft/Joly files.]}
}
{
}

\Verse{360}
{
    \zR{外面又报：“太医来了。”}
}
{
    \eR{[English source not present in the checked-in Bencraft/Joly files.]}
}
{
}

\Verse{361}
{
    \zR{贾琏接入，诊了一回。}
}
{
    \eR{[English source not present in the checked-in Bencraft/Joly files.]}
}
{
}

\Verse{362}
{
    \zR{大夫出来， 悄悄的告诉贾琏：“老太太的脉气不好，防着些。”}
}
{
    \eR{[English source not present in the checked-in Bencraft/Joly files.]}
}
{
}

\Verse{363}
{
    \zR{贾琏会意，与王夫人等说知。}
}
{
    \eR{[English source not present in the checked-in Bencraft/Joly files.]}
}
{
}

\Verse{364}
{
    \zR{王夫人即忙使眼色叫鸳鸯过来，叫他把老太太的装裹衣服预备出来。}
}
{
    \eR{[English source not present in the checked-in Bencraft/Joly files.]}
}
{
}

\Verse{365}
{
    \zR{鸳鸯自去料 理。}
}
{
    \eR{[English source not present in the checked-in Bencraft/Joly files.]}
}
{
}

\Verse{366}
{
    \zR{贾母睁眼要茶喝，邢夫人便进了一杯参汤。}
}
{
    \eR{[English source not present in the checked-in Bencraft/Joly files.]}
}
{
}

\Verse{367}
{
    \zR{贾母刚用嘴接着喝，便道：“不要 这个，倒一钟茶来我喝。”}
}
{
    \eR{[English source not present in the checked-in Bencraft/Joly files.]}
}
{
}

\Verse{368}
{
    \zR{众人不敢违拗，即忙送上来。}
}
{
    \eR{[English source not present in the checked-in Bencraft/Joly files.]}
}
{
}

\Verse{369}
{
    \zR{一口喝了，还要，又喝一 口，便说：“我要坐起来。”}
}
{
    \eR{[English source not present in the checked-in Bencraft/Joly files.]}
}
{
}

\Verse{370}
{
    \zR{贾政等道：“老太太要什么，只管说，可以不必坐起 来才好。”}
}
{
    \eR{[English source not present in the checked-in Bencraft/Joly files.]}
}
{
}

\Verse{371}
{
    \zR{贾母道：“我喝了口水，心里好些儿，略靠着和你们说说话儿。”}
}
{
    \eR{[English source not present in the checked-in Bencraft/Joly files.]}
}
{
}

\Verse{372}
{
    \zR{珍珠 等用手轻轻的扶起，看见贾母这会子精神好了些。}
}
{
    \eR{[English source not present in the checked-in Bencraft/Joly files.]}
}
{
}

\Verse{373}
{
    \zR{未知生死，下回分解。}
}
{
    \eR{[English source not present in the checked-in Bencraft/Joly files.]}
}
{
}

\Verse{374}
{
    \zR{(本章完)}
}
{
    \eR{[English source not present in the checked-in Bencraft/Joly files.]}
}
{
}
